\section{Putting the workers into Forge}
\label{sec:integration}

\subsection{Three new payloads, not a new authority model}

Forge's proposed runtime already distinguishes an ordinary candidate from a
conditional candidate whose premises still need proofs. The extension should
reuse that distinction and the existing acceptance boundary, rather than create
a second notion of accepted proof \cite{forge-runtime,forge-integration}.
The new data is the \emph{shape of the candidate}:
\begin{lstlisting}
AlgebraRegion:
  sourceFoldOccurrence
  selectedCarrierAndInstantiation
  localHoleTelescopes
  residualContracts
  sharedHoleConstraints
  reconstructionTerm

InductiveAlgebraCandidate:
  stateSchema
  initialSubstitution
  labeledTransitions
  equations
  actionMatricesOrPolynomialMultipliers
  targetCombination

RankedSCCCandidate:
  nodeTelescopes
  propositionFamilies
  branchProofConstructors
  typedCallSubstitutions
  rankAndDescentCertificates
\end{lstlisting}
These are proposed records. They are not currently wired into
\code{ForgeDesign.EngineReply}, and their fields do not themselves enforce
soundness. Their advantage is that every unresolved field corresponds to a
specific next action: prove a local contract, verify a polynomial identity,
transport a dependent argument, or discharge a descent condition.

\subsection{A target-directed worker protocol}

The polynomial worker receives an already reified transition problem and a
specific target equation. It may return a proved inductive family, an exact
encoded-model counterexample, a degree-escape report, or a bounded miss. The
first can become a conditional candidate for the original theorem once
initialization, transition semantics, and the target transport are supplied.
A degree-escape report can trigger the ideal lane; it cannot close the goal.

The continuation worker receives a fold occurrence and a residual specification.
It returns a typed term plus algebra-contract obligations. Successful local
search does not replace the final source elaboration. A term that fits the
neutral carrier but fails the source universe constraints is rejected, even if
its Python interpretation works on every test list.

The cyclic worker receives a candidate SCC and returns a ranked proof-program
skeleton. It does not emit ordinary checked facts for individual nodes before
the whole component has been lowered through well-founded induction. Its local
constructors can use algebraic facts from the invariant worker, but only after
those facts have been proved in the appropriate node context.

\subsection{A cooperative proof trace}

Consider a helper that performs either of the two polynomial updates from
Section~\ref{sec:spaces}, selected by a list of commands, and returns its final
state. The desired theorem is that the final state satisfies $z=2y$ when the
initial state is $(t,t^2,2t^2)$.

A direct induction on commands reaches two arithmetic step goals. Inducting only
on the target equation is insufficient: $z=2y$ alone does not imply that either
update preserves it for arbitrary $x,y,z$. The failed step is exactly the useful
signal to send to inductive algebra.

The space lane returns $e_1=e_2=0$ and the two constant action matrices.
Alternatively, target-seeded ideal completion begins with $z-2y$ and adds the
missing pullback equation. Forge strengthens the induction motive to the
resulting conjunction, proves the initialization by substitution, proves each
step by the action identities, and derives the original target as a linear
combination. No arbitrary induction predicate needs to be guessed beyond the
polynomial fragment.

Now split execution into two mutually recursive phases, with a transition
between phases that does not consume a command and a second transition that
consumes one. A plain ``list length decreases on every call'' gate fails. The
cyclic worker attaches phase-sensitive rank $2\,\operatorname{length}(xs)+b_v$,
then compiles the same strengthened motive through well-founded induction.
The action certificates solve arithmetic preservation; the rank certificates
solve recursion admissibility. These are different obligations with different
proof artifacts.

This combined end-to-end proof is an \emph{integration target}, not an executed
benchmark. The archive executes its algebra and descent components separately.
That separation is useful: it tells the implementer whether a later failure is
in mathematics, reification, dependent transport, or proof construction.

\subsection{Dependent contexts require explicit call morphisms}

A cyclic edge must map the callee's entire dependent telescope, not just the
visible arguments in its displayed goal. For example, if a node contains
$n:\N$, $xs:\operatorname{Vector}(A,n)$, and a hypothesis about $xs$, an edge
changing $n$ cannot keep the old vector or hypothesis by identity. It must supply
a vector at the new index and transport or reconstruct the dependent premise.

Similarly, a proof established under a hypothesis $H(n_0)$ about one fixed
current value cannot be silently generalized into a proof under $H(n)$ for all
$n$. The compiler should separate ambient parameters, which remain fixed for the
whole SCC, from state-dependent inputs and hypotheses, which are packaged in
$S_v$ or explicit premise arguments. Matching a back edge then constructs a
morphism of these packages and asks Lean to check it.

This is where a source-aware neutral representation earns its cost. Printed
variable names do not identify a telescope, and an alpha-normalized goal body
alone does not determine which hypotheses can be moved to a recursive call.
The necessary infrastructure is already part of Forge's identity design;
cyclic compilation adds a precise consumer that requires those identities to be
carried through a typed substitution.

\subsection{Reification obligations specific to these algorithms}

The polynomial algorithms operate over $\Q$ with simultaneous updates. A source
adapter must prove that its extracted map has those semantics. Sequential source
assignments need an SSA-style substitution pass before they become one map:
\code{x := x + 1; y := x + y} means $(x,y)\mapsto(x+1,x+y+1)$, not
$(x+1,x+y)$. This is an important negative integration fixture.

A natural-number subtraction is not unrestricted ring subtraction. It can be
translated into integer subtraction only under the appropriate comparison
premise, with that premise retained in the edge's guard and proved where the
edge is used. Division by a variable is not a polynomial update. Floating-point
arithmetic is not rational arithmetic. The first adapter should refuse these
unproved translations rather than send a nearby polynomial problem to a worker.

For a source recurrence over naturals, proving a polynomial identity after an
embedding into a characteristic-zero field can be useful, but the embedding and
reflection of the final equality are source-level obligations. A counterexample
with rational initialization parameters is not a counterexample to a theorem
quantified only over naturals unless those parameters are also shown to lie in
the source domain.

\subsection{Concrete module boundaries}

\begin{longtable}{@{}p{.40\textwidth}p{.55\textwidth}@{}}
\toprule
Proposed module & Responsibility and first acceptance condition \\
\midrule
\endhead
\code{Forge/Synthesis/AlgebraRegion.lean} & Extract one source fold occurrence, retain carrier/universe evidence, and construct local telescopes without losing binder identity. \\
\code{Forge/Synthesis/ContractJoin.lean} & Join shared-hole constraints; a correlated-hole fixture must reject incompatible projection products. \\
\code{Forge/Invariant/StateSchema.lean} & Reify a simultaneous rational-polynomial transition with a proved source interpretation. Sequential assignment and natural subtraction are mandatory controls. \\
\code{Forge/Invariant/PullbackSpace.lean} & Consume exact constant-action certificates; the coupled example must produce the original target theorem after kernel replay. \\
\code{Forge/Invariant/IdealReplay.lean} & Consume polynomial multipliers and exact orbit traces without invoking a CAS in the proof checker. \\
\code{Forge/Recursion/CallGraph.lean} & Construct typed node families and call substitutions, initially from explicit annotations. \\
\code{Forge/Recursion/RankReplay.lean} & Prove natural-domain and strict/lexicographic descent, including node phases. \\
\code{Forge/Recursion/CompileSCC.lean} & Lower all local constructors through one well-founded theorem; reject a missing branch even when the graph ranks successfully. \\
\bottomrule
\end{longtable}

The file names specify new responsibilities, not files claimed to exist in the
reviewed repository. A first implementation may combine modules, but it should
not combine their acceptance conditions.

\subsection{Library choices and substitutions}

\begin{table}[ht]
\small
\begin{tabularx}{\textwidth}{@{}p{.22\textwidth}YY@{}}
\toprule
Layer & Concrete choice & What still has to be implemented \\
\midrule
Exact matrix proposer & Python-FLINT \code{fmpq\_mat.rref}; integer nullspaces with \code{fmpz\_mat} & Preimage assembly, rational reconstruction where used, and action-certificate extraction. \\
Ideal proposer & Singular \code{lift}/\code{liftstd}; traced Buchberger for a minimal prototype & Original-generator provenance, substitutions, resource control, and multiplier serialization. \\
Small arithmetic holes & cvc5's SyGuS engine, optional & Source-scoped grammar and contract translation; candidate terms still need replay. \\
Lean polynomial semantics & Mathlib \code{MvPolynomial} evaluation and ring normalization & A reifier with interpretation lemmas and certificate-to-proof lowering. \\
Lean recursive proof & Core well-founded induction and recursive-definition elaboration & Tagged goal families, local proof constructors, source-level call transport. \\
\bottomrule
\end{tabularx}
\caption{Libraries are components, not finished adapters. References: \cite{flint-q,flint-z,singular,cvc5,mathlib-eval,lean-recursion}.}
\end{table}

The executable package does not require these proposer libraries. Its exact
search uses the Python standard library; SymPy 1.14.0 was used only for the
optional differential test. This keeps the evidence reproducible while making
the production replacement points explicit.

\subsection{Resource allocation that follows the mathematics}

The new workers should consume different budgets. The space worker is naturally
bounded by degree, monomial count, coefficient size, and elimination work. The
ideal worker needs generator, degree, basis-pair, coefficient, and process-memory
budgets. The synthesis worker needs local-candidate, join-row, whole-term, and
verification budgets. The cyclic worker needs graph-size, ranking-template,
solver, and proof-lowering budgets.

A useful escalation order is deterministic and explainable. Try a small algebra
region before enlarging the whole-term synthesis window. Try a low-degree space
before invoking ideal completion. Try scalar rankings with phases before
lexicographic tuples. Return the reason for escalation in the trace: an escaped
pullback degree, a failed local contract, or an SCC requiring more than one
ranking coordinate. These reasons are concrete mathematical diagnostics, not
an opaque score presented as evidence.

\section{Development sequence and release gates}
\label{sec:plan}

\subsection{First release: exact inductive algebra}

Implement \code{StateSchema} for explicit rational-polynomial recurrences and
lower the certificate identities to Lean proofs. Start with direct generated
\code{ring} proofs before investing in reflection. The entry theorem should
quantify over arbitrary transition words, not a sampled number of iterations.
Replay the coupled example using both the space and ideal certificates, and the
power-curve example using nonconstant multipliers.

The release gate is an ordinary theorem at the original target type, with
initialization and every transition label accounted for. Record proof construction
time, kernel checking time, and proof-term size separately. No successful
Python certificate should count toward this gate.

\subsection{Second release: one source-preserving algebra region}

Add continuation-local synthesis for the specific integer-indexing schema.
Preserve the exact source fold, default, integer domain, available primitives,
and dependent/universe evidence. Prove the derived local contract, then apply a
checked theorem for the source representation.

The original Leant indexing query is the acceptance target. Its isolated step
and the finite Python grammar are diagnostic fixtures only. First record the
original search branch's lexical identities, pending goals, and usage history;
then compare the same branch with and without region-local scheduling. This
isolates whether the new decomposition improves the actual miss rather than
merely making a different query easier.

\subsection{Third release: explicit cyclic skeletons}

Accept a user- or tactic-supplied cyclic skeleton with explicit branch constructors
and synthesize only its rank. Prove the two-phase, swapped-argument, and
lexicographic-growth examples in Lean. Reject the identity cycle, the alternating
replenishment cycle, and a graph with a missing logical base branch.

Only after these pass should Forge infer back edges from recurring obligations.
The first automatic mode should be limited to source families with straightforward
telescope transport. General dependent cyclic matching is a subsequent capability,
not a silent assumption of the initial compiler.

\subsection{What would constitute a meaningful improvement}

A production evaluation should have separate ablations for carrier selection,
local contracts, compatibility joins, invariant-space closure, ideal completion,
phase ranks, and lexicographic ranks. Each benchmark should retain its original
statement and source context. Measure solved goals under a fixed total budget,
not only candidate ordinals or search-only CPU time.

For mathematics-heavy recurrences, the key comparison is target-only induction
versus automatically discovered simultaneous invariants. For higher-order
construction, it is whole-term enumeration versus locally constrained assembly
under the same global accounting. For recursive proofs, it is fixed structural
induction versus a successfully lowered ranked SCC. A single aggregate score
would hide which new capability actually supplied the proof.
